\documentclass[a4paper]{article}
% Español (México) con XeLaTeX: fontspec para fuentes Unicode, polyglossia
% para la división silábica y los nombres del documento en la variante mexicana.
\usepackage{fontspec}
\usepackage{polyglossia}
\setdefaultlanguage[variant=mexican]{spanish}
% Los nombres chinos van como «pinyin (original)»: xeCJK da a los caracteres
% Han su propia fuente; sin ella XeLaTeX los omite con un aviso.
% FandolSong no cubre algunos caracteres de nombres (U+73FA); WenQuanYi Zen Hei sí.
\usepackage[AutoFallBack=true]{xeCJK}
\setCJKmainfont{FandolSong-Regular.otf}
\setCJKfallbackfamilyfont{\CJKrmdefault}{WenQuanYi Zen Hei}
% polyglossia no traduce los nombres de listings.
\AtBeginDocument{\providecommand{\lstlistingname}{}\renewcommand{\lstlistingname}{Listado}%
  \providecommand{\lstlistlistingname}{}\renewcommand{\lstlistlistingname}{Índice de listados}}
\usepackage{amsmath, amssymb}
\usepackage{graphicx}
\usepackage[margin=2.5cm]{geometry}
\usepackage[most]{tcolorbox}
\usepackage{etoolbox}
\usepackage{listings}
\usepackage{hyperref}
\usepackage{booktabs}
\usepackage{subcaption}
\usepackage{float}

\newtcolorbox{knowledgebox}[1]{enhanced, colback=blue!5!white, colframe=blue!75!black, colbacktitle=blue!75!black, coltitle=white, fonttitle=\bfseries, title={#1}, attach boxed title to top left={yshift=-2mm, xshift=2mm}, boxrule=1pt, sharp corners}
\newtcolorbox{importantbox}[1]{enhanced, colback=yellow!10!white, colframe=yellow!80!black, colbacktitle=yellow!80!black, coltitle=black, fonttitle=\bfseries, title={#1}, sharp corners}
\newtcolorbox{warningbox}[1]{enhanced, colback=red!5!white, colframe=red!75!black, colbacktitle=red!75!black, coltitle=white, fonttitle=\bfseries, title={#1}, sharp corners}

\lstset{language=Python, basicstyle=\ttfamily\small, keywordstyle=\color{blue}, stringstyle=\color{red!60!black}, commentstyle=\color{green!60!black}, breaklines=true, frame=single, numbers=left, numberstyle=\tiny\color{gray}, captionpos=b, extendedchars=true}

\newcommand{\notetitle}{Planning Agent: Plan-Execute-RePlan}
\newcommand{\noteauthors}{Elaboradas a partir de materiales públicos del curso}
\newcommand{\notedate}{2025}
\newcommand{\videochannel}{Wudaokou Nashi (五道口纳什)}
\newcommand{\videopublishdate}{2025}
\newcommand{\videoduration}{aprox. 14 minutos}
\newcommand{\videourl}{https://www.bilibili.com/video/BV1qa43z5EBJ}
\newcommand{\videocoverpath}{cover.jpg}
\newcommand{\repourl}{https://github.com/hqhq1025/ai-course-notes}


% Footer with repo info
\usepackage{fancyhdr}
\pagestyle{fancy}
\fancyhf{}
\fancyfoot[C]{\thepage}
\fancyfoot[R]{\footnotesize\color{gray}\href{\repourl}{AI Course Notes}}
\renewcommand{\headrulewidth}{0pt}
\renewcommand{\footrulewidth}{0pt}

\begin{document}
\begin{titlepage}
\centering
{\Large Notas del curso\par}\vspace{1.2cm}
{\huge\bfseries \notetitle\par}\vspace{0.8cm}
{\large \noteauthors\par}\vspace{0.3cm}
{\large \notedate\par}\vspace{1.2cm}
\ifdefempty{\videocoverpath}{}{\includegraphics[width=0.82\textwidth,height=0.45\textheight,keepaspectratio]{\videocoverpath}\par}
\vfill
\begin{tcolorbox}[width=0.9\textwidth, colback=black!2!white, colframe=black!60, sharp corners]
\textbf{Autor o canal del video}: \videochannel\par
\textbf{Fecha de publicación}: \videopublishdate\par
\textbf{Duración del video}: \videoduration\par
\ifdefempty{\videourl}{}{\textbf{Enlace del video}: \href{\videourl}{\nolinkurl{\videourl}}\par}\par
\textbf{Página del proyecto}: \href{\repourl}{\nolinkurl{\repourl}}\par
\end{tcolorbox}
\end{titlepage}
\tableofcontents
\newpage

\section{Panorama general de Plan-Execute-RePlan}

\subsection{Idea central}

En esta clase se presenta un Agentic Workflow muy general y ampliamente usado tanto en la investigación real como en el trabajo cotidiano: \textbf{Dynamic Plan-and-Solve}, es decir, el ciclo de tres pasos Plan-Execute-RePlan.

\begin{importantbox}{Las tres etapas de Plan-Execute-RePlan}
\begin{enumerate}
    \item \textbf{Plan}: al recibir la query/tarea compleja del usuario, primero se elabora una descomposición de sub-goals de arriba hacia abajo, produciendo una List of Steps (To-Do List)
    \item \textbf{Execute}: se procesa paso a paso cada step del Plan, obteniendo el outcome y el feedback mediante llamadas a herramientas
    \item \textbf{RePlan}: con base en los resultados de la ejecución y el feedback, se actualizan dinámicamente los steps restantes
\end{enumerate}
\end{importantbox}

Este patrón ya se ha adoptado ampliamente: Coding Agents como Cursor, Claude Code y Codex admiten un Plan Mode, en el que primero se elabora la To-Do List y después se ejecuta actualizándola sobre la marcha.

\subsection{Por qué funciona Plan-Execute-RePlan}

Un Workflow estructurado ayuda a:
\begin{itemize}
    \item estructurar el proceso de razonamiento del modelo
    \item reducir la complejidad global al resolver problemas complejos
    \item cerrar el ciclo mediante la descomposición de arriba hacia abajo y el feedback de abajo hacia arriba
\end{itemize}

\begin{knowledgebox}{Top-Down Plan vs. Bottom-Up Execute}
\begin{itemize}
    \item \textbf{Plan (Top-Down)}: descomposición de sub-goals de arriba hacia abajo, sin interacción real con el entorno; es un «cerebro en una cubeta», una secuencia de sub-goals puramente imaginada
    \item \textbf{Execute (Bottom-Up)}: interacción real con el entorno, obteniendo feedback real
    \item \textbf{RePlan}: con base en los resultados de la interacción real, se corrige dinámicamente el plan
\end{itemize}
\end{knowledgebox}

\subsection{Resumen de la sección}

Plan-Execute-RePlan es un Agentic Workflow muy general y validado en la práctica, cuyo valor central está en descomponer de forma estructurada los problemas complejos y, mediante un mecanismo de feedback dinámico, cerrar la brecha entre la planeación y la ejecución.
\section{Agentic Workflow vs. Autonomous Agent}

\subsection{Distinción de conceptos}

\begin{knowledgebox}{Definición de Agent según Anthropic}
De acuerdo con la definición de Agent que dio Anthropic a principios de año, Plan-Execute-RePlan pertenece al \textbf{Agentic Workflow predefinido}, no al Autonomous Agent. Pero esto no significa que el Agentic Workflow sea necesariamente peor que el Autonomous Agent: lo agentic es un espectro, y el Autonomous Agent es la forma final.
\end{knowledgebox}

La ventaja del Workflow predefinido está en ser \textbf{controlable, confiable y predecible}. Un Agent completamente autónomo (el paradigma ReAct puro) está limitado por la capacidad del modelo y puede tener un desempeño deficiente.

\subsection{Mejores prácticas de OpenAI sobre los Reasoning Model}

La posición oficial de la documentación de la API de OpenAI sobre los Reasoning Model:
\begin{itemize}
    \item \textbf{Modelos de razonamiento de la serie O}: adecuados para el papel de Planner, pueden dedicar más tiempo y un pensamiento más profundo a tareas complejas, formulando estrategias con mayor eficacia
    \item \textbf{Modelos no razonadores de GPT}: adecuados para el papel de Executor, buscan baja latencia y alta eficiencia de costo, ejecutando directamente tareas concretas
\end{itemize}

\subsection{Resumen de la sección}

El Agentic Workflow gana por ser controlable y confiable, y es la opción dominante en las aplicaciones prácticas actuales. La estrategia de usar un modelo fuerte para planear y uno débil para ejecutar coincide con la recomendación oficial de OpenAI.

\section{Implementación detallada en LangGraph}

\subsection{Definición del estado}

El estado (State) compartido globalmente contiene cuatro campos:
\begin{itemize}
    \item \texttt{input}: la pregunta o tarea del usuario
    \item \texttt{plan}: la Plan List formada (secuencia de sub-goal)
    \item \texttt{past\_steps}: los pasos ya ejecutados y sus resultados (de forma acumulativa)
    \item \texttt{response}: la respuesta final para el usuario
\end{itemize}

\subsection{Nodo Plan}

El diseño del System Prompt de Plan es muy general (independiente de la tarea concreta):
\begin{quote}
\textit{For the given objective, come up with a simple step by step plan. This plan should involve individual tasks, that if executed correctly will yield the correct answer. Do not add any superfluous steps. The result of the final step should be the final answer. Make sure that each step has all the information needed -- do not skip steps.}
\end{quote}

La salida es una lista estructurada de Steps.

\subsection{Nodo Execute}

Execute usa un ReAct Agent (\texttt{create\_react\_agent}) para procesar cada step del Plan. Un step puede corresponder a varias llamadas a herramientas. El resultado de la ejecución se agrega a \texttt{past\_steps}.

\subsection{Nodo RePlan}

La entrada de RePlan contiene tres partes:
\begin{enumerate}
    \item El objetivo original del usuario (objective)
    \item El Plan original
    \item Los steps ya ejecutados y sus resultados
\end{enumerate}

\begin{warningbox}{Restricciones clave del RePlan Prompt}
\begin{itemize}
    \item Si no hay más pasos por ejecutar, devolver directamente el control al usuario (return to user)
    \item De lo contrario, agregar al Plan solo los pasos que todavía deben ejecutarse
    \item \textbf{No devolver los steps ya ejecutados como parte del nuevo Plan}
\end{itemize}
\end{warningbox}

\subsection{Estructura del Graph}

Tres nodos + aristas condicionales:
\begin{enumerate}
    \item \textbf{Planner} $\rightarrow$ \textbf{Agent} (ReAct Execute) $\rightarrow$ \textbf{RePlan}
    \item RePlan tiene una arista condicional: decide si terminar (response) o continuar la ejecución (regresar a Agent)
\end{enumerate}

\subsection{Resumen de la sección}

La implementación en LangGraph muestra cómo llevar Plan-Execute-RePlan a un Graph claro de tres nodos, en el que Plan y RePlan determinan el límite superior de calidad del conjunto, mientras que Execute se encarga de las llamadas concretas a las herramientas.

\section{Economía y estrategias de combinación de modelos}

\begin{importantbox}{Principio de combinación de modelos fuertes y débiles}
\begin{itemize}
    \item \textbf{Planning y RePlan}: se usa un modelo fuerte (como un modelo de razonamiento), porque la planeación determina el límite superior de todo el Workflow
    \item \textbf{Execute}: se usa un modelo débil (como GPT-4.1 Nano/Mini), porque ya se le indicó al modelo qué hacer y solo necesita ejecutarlo paso a paso
\end{itemize}
Esto coincide por completo con el posicionamiento que hace OpenAI de los modelos de razonamiento y los que no lo son.
\end{importantbox}

Si la etapa de Planning es especialmente importante, también se puede introducir el patrón \textbf{Generator-Critique} para formar un Plan más sólido y confiable; no hay que esperar que el modelo proponga un buen Plan en un solo intento.

\subsection{Resumen de la sección}

Una estrategia razonable de combinación de modelos permite equilibrar el costo y el resultado: los modelos caros se encargan de la planeación y los modelos baratos, de la ejecución.

\section{Síntesis y ampliación}

\begin{enumerate}
    \item \textbf{Plan-Execute-RePlan es un Workflow muy general y eficaz}, aplicable a diversos escenarios de tareas complejas, y ha sido ampliamente validado.
    \item \textbf{Descomposición de arriba hacia abajo + retroalimentación de abajo hacia arriba}: Plan es la imaginación de un «cerebro en una cubeta», Execute es la interacción real, y RePlan cierra el gap entre ambos.
    \item \textbf{Consideraciones de economía}: un modelo fuerte hace la planeación y uno débil hace la ejecución, en concordancia con la recomendación oficial de OpenAI.
    \item \textbf{La calidad del Plan determina el límite superior}: si la complejidad de la tarea es alta, se puede usar un modelo más fuerte en la etapa de Planning o introducir un mecanismo de Generator-Critique.
    \item \textbf{De Agentic Workflow a Autonomous Agent}: la práctica actual demuestra que un Workflow predefinido supera a un Agent completamente autónomo en controlabilidad y confiabilidad.
\end{enumerate}

\end{document}
